% Source: https://artofproblemsolving.com/community/c3h1769463_rotation_and_reflections
% Stable identity: aops-topic-1769463-q01
Going round the circle, the three blue balloons split the nine red ones into three gaps of sizes $a,b,c\ge0$ with $a+b+c=9$. An arrangement is determined by the cyclic sequence of gaps, and rotations and reflections of the circle permute the three gaps cyclically or reverse their order. For three gaps these operations produce every ordering of $(a,b,c)$, so an arrangement corresponds exactly to an unordered triple $\{a,b,c\}$, that is, a partition of $9$ into at most $3$ parts:
\[
900,\ 810,\ 720,\ 711,\ 630,\ 621,\ 540,\ 531,\ 522,\ 441,\ 432,\ 333
\]
(written as $abc$ with $a\ge b\ge c$). There are $\boxed{12}$ arrangements.

Check with Burnside's lemma: the $24$ symmetries fix $\dbinom{12}{3}=220$ (identity), $4$ each (the two rotations by a third of a turn), $10$ each (the six reflections through two balloons: one fixed position blue and one symmetric pair blue) and $0$ for the rest, giving $\dfrac{220+8+60}{24}=12$.
